% GENERATED FILE. Edit canonical lessons, then run npm run generate:books.
% canonical-source-hash: fnv1a64:d2d5c0266d8ec83d
% canonical-lessons: TA-C170-ovvamai, TA-C170-torru, TA-C170-talumpu, TA-C170-vayirruppokku, TA-C170-mulai

\chapter{Allergy, Infection, Scar, Diarrhoea, Brain}
\label{ch:ta-allergy-infection-scar-diarrhoea-brain}
\hlchaptermodality{170}

\begin{chapteropening}
\textbf{By the end of this chapter:} \emph{I can say an allergy, an infection, a scar, diarrhoea and the brain.}

Complete the last lesson of chapter 170: I can say an allergy, an infection, a scar, diarrhoea and the brain.
\end{chapteropening}

\section[\texorpdfstring{\ta{ஒவ்வாமை} (ovvāmai)}{ஒவ்வாமை (ovvāmai)}]{\ta{ஒவ்வாமை} (ovvāmai) — an allergy}

\label{lesson:TA-C170-ovvamai}



\begin{quote}
Before the new one: say the Tamil for the heart, then the Tamil for blood.
\end{quote}

\subsection*{You'll want to know: \ta{ஒவ்வாமை}}
\textbf{\ta{ஒவ்வாமை}} — \emph{ovvāmai} — \textquotedblleft{}an allergy\textquotedblright{}.

\begin{grammarlens}[title={What you've built}]
One more word for this chapter.
\end{grammarlens}

\subsection*{Your turn}
\begin{itemize}
\raggedright
  \item \emph{Say it:} \emph{ovvāmai}
  \item \emph{Say it:} \emph{ovvāmai}, once more
  \item \emph{Say it:} say \emph{ovvāmai}
  \item \emph{Recall:} say the Tamil for the heart, then the Tamil for blood, then say \emph{ovvāmai} again
\end{itemize}

\begin{tcolorbox}[breakable,colback=teal!4,colframe=teal!35!black,title={Before you move on}]
What does \ta{ஒவ்வாமை} mean? (\textquotedblleft{}An allergy\textquotedblright{}.) Say it once more.
\end{tcolorbox}
\section[\texorpdfstring{\ta{தொற்று} (toṟṟu)}{தொற்று (toṟṟu)}]{\ta{தொற்று} (toṟṟu) — an infection}

\label{lesson:TA-C170-torru}



\begin{quote}
Before the new one: say the Tamil for blood, then the Tamil for an allergy.
\end{quote}

\subsection*{You'll want to know: \ta{தொற்று}}
\textbf{\ta{தொற்று}} — \emph{toṟṟu} — \textquotedblleft{}an infection\textquotedblright{}.

\begin{grammarlens}[title={What you've built}]
One more word for this chapter.
\end{grammarlens}

\subsection*{Your turn}
\begin{itemize}
\raggedright
  \item \emph{Say it:} \emph{toṟṟu}
  \item \emph{Say it:} \emph{toṟṟu}, once more
  \item \emph{Say it:} say \emph{toṟṟu}
  \item \emph{Recall:} say the Tamil for blood, then the Tamil for an allergy, then say \emph{toṟṟu} again
\end{itemize}

\begin{tcolorbox}[breakable,colback=teal!4,colframe=teal!35!black,title={Before you move on}]
What does \ta{தொற்று} mean? (\textquotedblleft{}An infection\textquotedblright{}.) Say it once more.
\end{tcolorbox}
\section[\texorpdfstring{\ta{தழும்பு} (taḻumpu)}{தழும்பு (taḻumpu)}]{\ta{தழும்பு} (taḻumpu) — a scar}

\label{lesson:TA-C170-talumpu}



\begin{quote}
Before the new one: say the Tamil for an allergy, then the Tamil for an infection.
\end{quote}

\subsection*{You'll want to know: \ta{தழும்பு}}
\textbf{\ta{தழும்பு}} — \emph{taḻumpu} — \textquotedblleft{}a scar\textquotedblright{}.

\begin{grammarlens}[title={What you've built}]
One more word for this chapter.
\end{grammarlens}

\subsection*{Your turn}
\begin{itemize}
\raggedright
  \item \emph{Say it:} \emph{taḻumpu}
  \item \emph{Say it:} \emph{taḻumpu}, once more
  \item \emph{Say it:} say \emph{taḻumpu}
  \item \emph{Recall:} say the Tamil for an allergy, then the Tamil for an infection, then say \emph{taḻumpu} again
\end{itemize}

\begin{tcolorbox}[breakable,colback=teal!4,colframe=teal!35!black,title={Before you move on}]
What does \ta{தழும்பு} mean? (\textquotedblleft{}A scar\textquotedblright{}.) Say it once more.
\end{tcolorbox}
\section[\texorpdfstring{\ta{வயிற்றுப்போக்கு} (vayiṟṟuppōkku)}{வயிற்றுப்போக்கு (vayiṟṟuppōkku)}]{\ta{வயிற்றுப்போக்கு} (vayiṟṟuppōkku) — diarrhoea}

\label{lesson:TA-C170-vayirruppokku}



\begin{quote}
Before the new one: say the Tamil for an infection, then the Tamil for a scar.
\end{quote}

\subsection*{You'll want to know: \ta{வயிற்றுப்போக்கு}}
\textbf{\ta{வயிற்றுப்போக்கு}} — \emph{vayiṟṟuppōkku} — \textquotedblleft{}diarrhoea\textquotedblright{}.

\begin{grammarlens}[title={What you've built}]
One more word for this chapter.
\end{grammarlens}

\subsection*{Your turn}
\begin{itemize}
\raggedright
  \item \emph{Say it:} \emph{vayiṟṟuppōkku}
  \item \emph{Say it:} \emph{vayiṟṟuppōkku}, once more
  \item \emph{Say it:} say \emph{vayiṟṟuppōkku}
  \item \emph{Recall:} say the Tamil for an infection, then the Tamil for a scar, then say \emph{vayiṟṟuppōkku} again
\end{itemize}

\begin{tcolorbox}[breakable,colback=teal!4,colframe=teal!35!black,title={Before you move on}]
What does \ta{வயிற்றுப்போக்கு} mean? (\textquotedblleft{}Diarrhoea\textquotedblright{}.) Say it once more.
\end{tcolorbox}
\section[\texorpdfstring{\ta{மூளை} (mūḷai)}{மூளை (mūḷai)}]{\ta{மூளை} (mūḷai) — the brain}

\label{lesson:TA-C170-mulai}



\begin{quote}
Before the new one: say the Tamil for a scar, then the Tamil for diarrhoea.
\end{quote}

\subsection*{You'll want to know: \ta{மூளை}}
\textbf{\ta{மூளை}} — \emph{mūḷai} — \textquotedblleft{}the brain\textquotedblright{}.

\begin{grammarlens}[title={What you've built}]
One more word for this chapter.
\end{grammarlens}

\subsection*{Your turn}
\begin{itemize}
\raggedright
  \item \emph{Say it:} \emph{mūḷai}
  \item \emph{Say it:} \emph{mūḷai}, once more
  \item \emph{Say it:} say \emph{mūḷai}
  \item \emph{Recall:} say the Tamil for a scar, then the Tamil for diarrhoea, then say \emph{mūḷai} again
\end{itemize}

\begin{tcolorbox}[breakable,colback=teal!4,colframe=teal!35!black,title={Before you move on}]
What does \ta{மூளை} mean? (\textquotedblleft{}The brain\textquotedblright{}.) Say it once more.
\end{tcolorbox}
